\begin{table}[!t]
\centering\footnotesize
\setlength{\tabcolsep}{4pt}
\caption{Endpoint Ablations With Control and Sensitivity Rows}\label{tab:ablation_full}
\begin{adjustbox}{max width=\columnwidth}
\begin{tabular}{lrrrr}
\toprule
Variant & Cells & Mean & Min. & Max. \\
\midrule
Fed-MDBSCAN-G, stage 1 only & 12 & $-$0.081 & $-$1.50 & $+$0.84 \\
Geometric-median distance control & 12 & $-$0.081 & $-$1.50 & $+$0.84 \\
Fed-MDBSCAN-G, no activation memory & 12 & 0.000 & 0.00 & 0.00 \\
Fed-MDBSCAN-G, no valve & 12 & $+$0.278 & $-$0.27 & $+$3.03 \\
Trust-region multiplier 1.5 & 12 & $-$0.916 & $-$8.42 & $+$2.90 \\
Trust-region multiplier 2.0 & 12 & $-$0.860 & $-$3.09 & $+$2.95 \\
Trust-region multiplier 3.0 & 12 & $+$0.575 & $-$6.38 & $+$12.82 \\
\bottomrule
\end{tabular}
\end{adjustbox}

\par\vspace{3pt}\parbox{\linewidth}{\footnotesize Variant minus full composite, final accuracy in percentage points; cells are dataset~$\times$~condition means over three seeds in the matched ablation block. Ranges describe heterogeneity across conditions, not confidence intervals. The geometric-median distance control is the standalone Stage-1 rule, a FedG2L filter with threshold factor 2.5; its accuracy, FPR, and TPR equal those of the stage-1-only variant in all 36 runs. The trust-region rows set the Stage-1 radius factor of the full composite to 1.5, 2.0, and 3.0 instead of 2.5.}
\end{table}
